\section{12. Dashboard}

The dashboard is one process, \code{viz}, with two faces. It builds both from
what the control system itself sees: the observations, the decisions and the
actuator events on MQTT. The true state of the rooms is not on any topic it may
read, so the dashboard shows what was measured, never what was true.

\subsection{12.1 The building now: layers and alerts in the BuildSim viewer}

viz writes five room layers to BuildSim once per real second, and the 3D viewer
colours each room by the layer chosen. The viewer itself is provided by the
course; what is ours is what is mapped into it.

\begin{longtable}[]{@{}L{0.24\linewidth}L{0.18\linewidth}L{0.16\linewidth}L{0.32\linewidth}@{}}
\toprule
\textbf{Layer} & \textbf{Source} & \textbf{Range} & \textbf{Colours} \\
\midrule
\endhead
\COtwo{} & sensor reading & 400--1600 ppm & Blue to red around grey at the 1000~ppm limit \\
People counted & sensor reading & 0--50 & Light to dark blue \\
People expected soon & prediction & 0--50 & Light to dark blue \\
Airflow & actuator state & 0--100~\% of the room's maximum & Light to dark blue \\
Temperature & sensor reading & 17--27 °C & Blue to red around grey at 22~°C \\
\bottomrule
\end{longtable}

A quantity of ``how much'' is one hue from light to dark; a quantity on either
side of a target is blue and red around a neutral grey that sits exactly on the
target. Airflow is shown as a share of each room's own maximum, because an
office's 20~l/s and a lecture room's 900~l/s do not fit one scale. A room drops
out of a layer when its reading is more than 180~s old or its sensor has been
found stuck: either colour would claim something nobody measured.

Alerts are the current picture, at most 100, most severe first:

\begin{longtable}[]{@{}L{0.11\linewidth}L{0.36\linewidth}L{0.22\linewidth}L{0.21\linewidth}@{}}
\toprule
\textbf{Severity} & \textbf{Alert} & \textbf{Raised by} & \textbf{Cleared by} \\
\midrule
\endhead
Critical & \COtwo{} override in a room & event \code{override\_co2} & \code{override\_released} \\
Warning & No reading from a room's sensor & a reading older than 180~s & its next reading \\
Warning & A sensor looks stuck, with what was seen and what the system does about it & a decision with a \code{fault} & the next decision without one \\
Warning & Damper or radiator valve on fallback & event \code{ttl\_expired} & the next accepted command \\
Warning & Command rejected & event \code{rejected} & five minutes later \\
Info & No trained model, once for the building & decisions in mode \code{no\_forecast} & normal decisions \\
\bottomrule
\end{longtable}

\missingfigure{Screenshot of the BuildSim viewer with the \COtwo{} layer and at
least one alert, for example during \code{deploy/fault-tests.sh}:
\code{http://127.0.0.1:9090/?layer=co2\&floor=level0}.}

\subsection{12.2 The building over time: the graphs page}

The viewer answers ``what is the building doing now''. Most of what a
ventilation system is asked is about time: did the air stay good through the
lecture, when did the damper open, what did the fault do. viz therefore serves
its own page, one HTML file with no outside dependencies, so it also works
without internet.

\begin{itemize}
\tightlist
\item \textbf{The building:} people counted, ventilation power and its share
  of what constant flow would use, rooms above the \COtwo{} limit, and alerts;
  three graphs over time; the alert log; a table of every room.
\item \textbf{One room:} \COtwo{} with the limit and the override level; people
  counted and people ventilated for; airflow commanded and reached, with design
  and maximum flow; temperature and setpoint with the comfort band. The times
  an alert was active are shaded in all four.
\item The range is the last hour, 6 hours, 24 hours or 7 days. Pointing at a
  moment in one graph shows it in all of them, by mouse or keyboard, and every
  graph has its values as a table underneath.
\end{itemize}

In every graph blue is what was measured or reached and orange is what the
controller expected or commanded; yellow and red mean an alert and nothing
else. A minute in which nothing was seen is a hole in the line, never a zero
and never the previous value. The one derived number, ventilation power, is
the summed reached airflow times the energy proxy's constants for fans and
supply-air heating; the radiators are left out because their power depends on
the true room temperature, which the dashboard never sees.

The history is held in viz's memory, one point per room and model minute for a
week. Storage remains the owner of the history: after a restart viz reads the
run so far from it, so a killed dashboard comes back with its graphs (T-39).

\missingfigure{Screenshot of the graphs page, room view, showing a fault: for
example room 1570 during the stuck \COtwo{} sensor of T-36 (flat \COtwo{} line,
shaded alert period, the damper's step to design flow).
\code{http://127.0.0.1:8060}.}

What the graphs showed that the tables had not: summed over the building, the
predictive policy's ``people ventilated for'' runs well ahead of the people
counted in the morning (57 against 6 at 09:29 on a development run). That gap
is the extra energy over the reactive policy, as a picture.
